\documentclass[11pt,a4paper]{extarticle}
% Shared preamble for blueMarine LaTeX reports
% % Shared preamble for blueMarine LaTeX reports
% % Shared preamble for blueMarine LaTeX reports
% \input{preamble.tex} after \documentclass{...}

\usepackage[T1]{fontenc}
\usepackage[utf8]{inputenc}
\usepackage{lmodern}

\usepackage[margin=2.5cm]{geometry}
\setlength{\parskip}{0.6em}
\setlength{\parindent}{0pt}

\usepackage{amsmath,amssymb}
\usepackage{booktabs}
\usepackage{array}
\usepackage{tabularx}
\usepackage{longtable}
\usepackage[table,dvipsnames]{xcolor}
\usepackage{colortbl}
\usepackage{graphicx}
\usepackage{float}
\usepackage{placeins} % \FloatBarrier

\usepackage[
  colorlinks=true,
  linkcolor=MidnightBlue,
  citecolor=MidnightBlue,
  urlcolor=MidnightBlue
]{hyperref}

\usepackage{enumitem}
\setlist[itemize]{topsep=0.3em,itemsep=0.2em}
\setlist[enumerate]{topsep=0.3em,itemsep=0.2em}

% Optional: comment out if draft watermarks not wanted
% \usepackage{draftwatermark}
% \SetWatermarkText{DRAFT}
% \SetWatermarkScale{0.8}

% Convenience macros
\newcommand{\Fmsy}{F_{\mathrm{MSY}}}
\newcommand{\Blim}{B_{\lim}}
\newcommand{\Btrig}{MSYB_{\mathrm{trigger}}}
\newcommand{\TODO}[1]{\textcolor{red}{\textbf{[TODO: #1]}}}
\newcommand{\figpath}{figs/}

% Placeholder figure. Arg 2 = file path (underscores allowed via detokenize).
\newcommand{\placeholderfig}[2]{%
  \begin{figure}[H]
  \centering
  \fbox{\parbox{0.9\textwidth}{\centering\vspace{3cm}
    \textit{Figure placeholder: #1}\\[0.5em]
    {\ttfamily\detokenize{#2}}\vspace{3cm}}}
  \caption{#1}
  \end{figure}%
}
 after \documentclass{...}

\usepackage[T1]{fontenc}
\usepackage[utf8]{inputenc}
\usepackage{lmodern}

\usepackage[margin=2.5cm]{geometry}
\setlength{\parskip}{0.6em}
\setlength{\parindent}{0pt}

\usepackage{amsmath,amssymb}
\usepackage{booktabs}
\usepackage{array}
\usepackage{tabularx}
\usepackage{longtable}
\usepackage[table,dvipsnames]{xcolor}
\usepackage{colortbl}
\usepackage{graphicx}
\usepackage{float}
\usepackage{placeins} % \FloatBarrier

\usepackage[
  colorlinks=true,
  linkcolor=MidnightBlue,
  citecolor=MidnightBlue,
  urlcolor=MidnightBlue
]{hyperref}

\usepackage{enumitem}
\setlist[itemize]{topsep=0.3em,itemsep=0.2em}
\setlist[enumerate]{topsep=0.3em,itemsep=0.2em}

% Optional: comment out if draft watermarks not wanted
% \usepackage{draftwatermark}
% \SetWatermarkText{DRAFT}
% \SetWatermarkScale{0.8}

% Convenience macros
\newcommand{\Fmsy}{F_{\mathrm{MSY}}}
\newcommand{\Blim}{B_{\lim}}
\newcommand{\Btrig}{MSYB_{\mathrm{trigger}}}
\newcommand{\TODO}[1]{\textcolor{red}{\textbf{[TODO: #1]}}}
\newcommand{\figpath}{figs/}

% Placeholder figure. Arg 2 = file path (underscores allowed via detokenize).
\newcommand{\placeholderfig}[2]{%
  \begin{figure}[H]
  \centering
  \fbox{\parbox{0.9\textwidth}{\centering\vspace{3cm}
    \textit{Figure placeholder: #1}\\[0.5em]
    {\ttfamily\detokenize{#2}}\vspace{3cm}}}
  \caption{#1}
  \end{figure}%
}
 after \documentclass{...}

\usepackage[T1]{fontenc}
\usepackage[utf8]{inputenc}
\usepackage{lmodern}

\usepackage[margin=2.5cm]{geometry}
\setlength{\parskip}{0.6em}
\setlength{\parindent}{0pt}

\usepackage{amsmath,amssymb}
\usepackage{booktabs}
\usepackage{array}
\usepackage{tabularx}
\usepackage{longtable}
\usepackage[table,dvipsnames]{xcolor}
\usepackage{colortbl}
\usepackage{graphicx}
\usepackage{float}
\usepackage{placeins} % \FloatBarrier

\usepackage[
  colorlinks=true,
  linkcolor=MidnightBlue,
  citecolor=MidnightBlue,
  urlcolor=MidnightBlue
]{hyperref}

\usepackage{enumitem}
\setlist[itemize]{topsep=0.3em,itemsep=0.2em}
\setlist[enumerate]{topsep=0.3em,itemsep=0.2em}

% Optional: comment out if draft watermarks not wanted
% \usepackage{draftwatermark}
% \SetWatermarkText{DRAFT}
% \SetWatermarkScale{0.8}

% Convenience macros
\newcommand{\Fmsy}{F_{\mathrm{MSY}}}
\newcommand{\Blim}{B_{\lim}}
\newcommand{\Btrig}{MSYB_{\mathrm{trigger}}}
\newcommand{\TODO}[1]{\textcolor{red}{\textbf{[TODO: #1]}}}
\newcommand{\figpath}{figs/}

% Placeholder figure. Arg 2 = file path (underscores allowed via detokenize).
\newcommand{\placeholderfig}[2]{%
  \begin{figure}[H]
  \centering
  \fbox{\parbox{0.9\textwidth}{\centering\vspace{3cm}
    \textit{Figure placeholder: #1}\\[0.5em]
    {\ttfamily\detokenize{#2}}\vspace{3cm}}}
  \caption{#1}
  \end{figure}%
}


%==============================================================================
% Technical report template — Backtest of ICES Advice
%
% Compile from tex/:
%   pdflatex report_template.tex
%
% Week~1 content can be merged from week1_summary.tex / generated/*.tex.
% Analysis notebooks write data/results/*.RData; figures under tex/figs/.
%==============================================================================

\title{Backtest of ICES Advice for Selected Northeast Atlantic Stocks\\
\large Technical report}
\author{%
  Laurence Kell\\
  \small Sea++
}
\date{\small \today}

\begin{document}
\maketitle

\begin{abstract}
\noindent
\TODO{2--3 paragraph abstract summarising aims, case studies, methods
(perfect-management peels), and headline findings once backtests are complete.}
\end{abstract}

\tableofcontents
\newpage

%------------------------------------------------------------------------------
\section{Introduction}
\label{sec:intro}

\TODO{Motivation: need to separate assessment/advice robustness from
implementation of advice; policy relevance for HCRs and rebuilding.}

This report documents a backtest of fisheries advice for selected Northeast
Atlantic stocks. Stock assessments are reconstructed annually over a historical
period (retrospective peels). The advice rule in operation each year is applied,
and stocks are projected forward under perfect implementation of catch advice.
Counterfactual outcomes are compared with realised catches, stock status, and
risks associated with overfishing.

\subsection{Aims and objectives}

\begin{enumerate}
  \item \textbf{Reconstruct} the annual perception of stock status using only
        data and methods available at the time.
  \item \textbf{Reproduce} the harvest control rule (ICES advice rule) in force
        each year.
  \item \textbf{Simulate} forward under perfect management (catches match
        advised catch).
  \item \textbf{Compare} counterfactual outcomes with realised history.
\end{enumerate}

\subsection{Report structure}

Section~\ref{sec:methods} describes the backtest procedure, including
operating-model conditioning (Section~\ref{sec:om-condition}).
Section~\ref{sec:casestudies} lists confirmed case-study stocks.
Section~\ref{sec:screening} summarises week~1 screening and benchmarks.
Section~\ref{sec:results} presents stock-specific results.
Section~\ref{sec:synthesis} provides a cross-stock synthesis.
Section~\ref{sec:discussion} discusses implications and limitations.
Section~\ref{sec:software} documents software and reproducibility.

%------------------------------------------------------------------------------
\section{Methods: backtest procedure}
\label{sec:methods}

% Generic peer-review Methods (OM, observation models, HCR, metrics, software):
% \input{C:/active/flr/backtest/inst/methods/methods-backtest.tex}
% Application choices for this study are summarised below; full generic text
% is in FLBacktest vignette methods-backtest / inst/methods/methods-backtest.tex.

The study separates two questions that are often conflated: whether the
assessment and advice rule are robust, and whether management outcomes differ
because advice was not followed. The closed-loop procedure (perceive--advise--
implement--compare) follows the \texttt{FLBacktest} Methods vignette. Historical
outcomes may differ from those expected under advice not only because recommended
catches were not implemented exactly, but also because stock perception,
reference points, and recruitment changed through time. The backtest separates
these effects for the six Northeast Atlantic ICES Category~1 stocks in
Table~\ref{tab:stocks}.

\subsection{Historical assessment reconstruction}

\TODO{Detail peel configurations (SAM / equivalent), data cuts, validation
against published ICES assessments, and any shortcuts used.}

For each stock and peel year \(t\), the assessment is re-run or reconstructed
using only contemporaneous information. The peel supplies the perception of
SSB, \(F\), and recruitment used in the advice step---not the latest revised
history. The effective window is about 7--10 assessment years (roughly
2016--2026), set by SAG online coverage. Short-cut runs may replace full peels
with multiplicative SSB error conditioned on retrospective bias
(\texttt{retroBias}, \texttt{shortcut}). This is a simulation experiment, not a
statutory re-assessment.

\subsection{Advice rule}

\begin{itemize}
  \item \textbf{Celtic and Irish Sea gadoids:} ICES Category~1 MSY approach---
        \(F = \Fmsy\) above \(\Btrig\), reduced linearly towards \(\Blim\),
        and zero at or below \(\Blim\).
  \item \textbf{Northeast Atlantic mackerel:} Coastal States / ICES MSY rule
        (including any stability constraints applying in year \(t\)).
\end{itemize}
Reference points are those published or used for advice in year \(t\).
Benchmark years are treated as structural breaks when \(\Fmsy\), \(\Blim\),
or \(\Btrig\) are redefined.

\subsection{Operating-model conditioning}
\label{sec:om-condition}

The OM representing \emph{true} stock dynamics for each case study is taken
from the most recent accepted age-based assessment, stored as a local
\texttt{FLStock}. The six OMs cover Celtic Sea cod, whiting, and haddock;
Irish Sea cod and whiting; and Northeast Atlantic mackerel. Conditioning prepares
each assessment for closed-loop
projection under \texttt{FLBacktest::hcrICES} and \texttt{FLasher}. The
critical biological step is the \emph{stock--recruitment relationship}: without
an assessment-consistent relationship and recruitment deviates, short-term
counterfactuals conflate fishing-policy effects with year-class strength
(Section~\ref{sec:screening}).

\paragraph{Data.}
Each OM uses the windowed \texttt{FLStock} (SSB, \(F\), recruitment, catch,
biology, and selectivity) from \texttt{data/reference/stocks.csv}; SAG
assessment-year series and ICES advice over the peel window; and HCR
parameters \(\Fmsy\), \(\Blim\), and \(\Btrig\) from SAG (\texttt{FLBRP} or
recent assessment quantities only as documented fallbacks).

\paragraph{Procedure.}
For each stock:
\begin{enumerate}
  \item Load and window the \texttt{FLStock} to the advice year; clean discard
        and landings slots so that \texttt{FLasher} projections are well
        defined.
  \item Fit a Beverton--Holt stock--recruitment relationship with
        \texttt{icesdata::eql} (\texttt{ftmb2}, scalar SPR0, steepness and
        \(R_0\) priors). Year-varying SPR0 diagnostics use
        \texttt{icesdata::ftmb} / \texttt{eqlFn}. Dynamic \(B_0\) without an
        SRR is \(R_t \times \mathrm{SPR}_{0,t}\). Year-specific \(v_t\) TMB
        fits (\texttt{FLRebuild::ftmb\_b0dyn}) are not used for the production
        OM. Unfished biomass can still vary with historical biology and
        selectivity in the diagnostic series. Retain the fitted relationship
        and log-scale recruitment deviates.
  \item Build an \texttt{FLBRP} equilibrium object and attach HCR parameters
        \texttt{ftar}, \texttt{blim}, and \texttt{btrig} from SAG. An
        implausible FLBRP \(\Fmsy\) is not substituted for the SAG advice
        parameter.
\end{enumerate}
Conditioned OMs are written to \texttt{data/results/03.0\_om.RData}
(\texttt{Rmd/02.0\_condition\_om.Rmd}).

\paragraph{Recruitment treatments.}
Screening used geometric-mean recruitment as a simple proxy to show how
strongly trajectories depend on year-class strength; that proxy is not the
recruitment model for the full backtest. Counterfactual runs use the fitted
stock--recruitment relationship and propagate recruitment deviates. Deterministic
mean-recruitment runs are retained only as a sensitivity comparison, so that
fishing mortality and recruitment can be separated: comparing mean vs deviate
paths holds policy fixed and isolates year-class effects; comparing deviate-informed
paths with the realised assessment isolates departures of realised catch (or
\(F\)) from advice.

\paragraph{Case-study caveats.}
Northeast Atlantic mackerel requires care at the recent benchmark (scale and
reference-point revision); results remain provisional until benchmark-year
handling is agreed. Irish Sea whiting has sparse recent SAG coverage and was
windowed to an earlier terminal year; it remains in the set with peel
diagnostics flagged.

SAG availability limits the stock-specific retrospective peel windows to
approximately 7--10 assessment years. This is a simulation study, not a
statutory re-assessment of any of the six stocks.

\subsection{Perfect-management projection}

Closed-loop simulation uses \texttt{FLBacktest::hcrICES} with
\texttt{FLasher} projections. From each peel (or short-cut perception):
\begin{enumerate}
  \item Reconstruct the stock status perceived from information available in
        year~\(t\) (not the latest revised estimate).
  \item Apply the contemporaneous HCR and its reference points to obtain
        advised catch for the next management year.
  \item Set realised catch equal to advised catch (perfect implementation: no
        TAC overshoot, override, or implementation error) and project the
        conditioned OM forward under the chosen recruitment treatment.
\end{enumerate}

\paragraph{Conceptual contrast.}
Over the backtest window we compare:
\begin{enumerate}
  \item \textbf{Realised trajectory:} history implied by the conditioned OM
        when forced by observed catch (or realised \(F\)), with assessed
        recruitment deviates.
  \item \textbf{Counterfactual trajectory:} history under perfect
        implementation of contemporaneous advice in the same OM.
\end{enumerate}
Both use the same conditioned biological dynamics and stock--recruitment
relationship. Differences describe the effect of the counterfactual fishing
history under the stated OM assumptions; they are not evidence that a new
assessment should replace the accepted ICES assessment. The short peel window
and OM assumptions bound the inferences.

\subsection{Evaluation metrics}

\begin{itemize}
  \item \textbf{Biological status} --- time below \(\Blim\) and \(\Btrig\);
        rebuilding times to \(\Btrig\) or another specified threshold;
  \item \textbf{Exploitation} --- realised and counterfactual \(F\) relative
        to \(\Fmsy\);
  \item \textbf{Yield and stability} --- annual and cumulative catch, and
        interannual variability.
\end{itemize}

%------------------------------------------------------------------------------
\section{Case-study stocks}
\label{sec:casestudies}

\TODO{Replace with generated/week1\_tab\_stocks.tex or an updated table once the list is final.}

All six proposal candidates were retained after week~1 screening
(see Section~\ref{sec:screening}).

\begin{table}[H]
\centering
\caption{Confirmed case-study stocks.}
\label{tab:stocks}
\begin{tabular}{@{}llll@{}}
\toprule
Stock & SID & ICES area & Advice basis \\
\midrule
Celtic Sea cod & \texttt{cod.27.7e-k} & 7.e--k & ICES MSY approach \\
Celtic Sea whiting & \texttt{whg.27.7b-ce-k} & 7.b--c, 7.e--k & ICES MSY approach \\
Celtic Sea haddock & \texttt{had.27.7b-k} & 7.b--k & ICES MSY approach \\
Irish Sea whiting & \texttt{whg.27.7a} & 7.a & ICES MSY approach \\
Irish Sea cod & \texttt{cod.27.7a} & 7.a & ICES MSY approach \\
NE Atlantic mackerel & \texttt{mac.27.nea} & 1--8, 14, 9.a & MSY rule (Coastal States) \\
\bottomrule
\end{tabular}
\end{table}

%------------------------------------------------------------------------------
\section{Week~1 screening and benchmarks}
\label{sec:screening}

\TODO{Optionally \texttt{\textbackslash input} narrative or tables from
\texttt{week1\_summary.tex} / \texttt{generated/}. For a short report, keep a
condensed summary and cite the standalone week~1 PDF.}

Week~1 confirmed case studies, screened SAG coverage and latest status, and
ran historical projections at target \(F\). Key discussion points for the
full backtest:
\begin{itemize}
  \item Perception of the stock changes at each update or benchmark.
  \item Recruitment has a large impact on short-term trajectories.
  \item Reference points are routinely changed at a benchmark.
\end{itemize}

\placeholderfig{SAG SSB retrospectives (week 1)}{figs/week1_sag_ssb.pdf}
\placeholderfig{SAG vs FLStock scaled trajectories (week 1)}{figs/week1_sag_fl.pdf}
\placeholderfig{Historical projections at target F (week 1)}{figs/week1_projections.pdf}

% Prefer live figures when files exist (uncomment and remove placeholders):
% \begin{figure}[H]\centering
% \includegraphics[width=\textwidth]{\figpath week1_sag_ssb.pdf}
% \caption{SAG SSB by assessment year.}\label{fig:sag-ssb}\end{figure}
%------------------------------------------------------------------------------
\section{Results by stock}
\label{sec:results}

\TODO{One subsection per stock. Pattern below; duplicate and fill.}

\subsection{Celtic Sea cod (\texttt{cod.27.7e-k})}
\label{sec:cod.27.7e-k}

\paragraph{Assessment reconstruction.}
\TODO{Peel years, diagnostics, agreement with published assessments.}

\paragraph{Advice and perfect-management counterfactual.}
\TODO{Describe HCR application and projection settings.}

\placeholderfig{Celtic Sea cod: realised vs counterfactual SSB}{figs/results_cod27.7e-k_ssb.pdf}

\paragraph{Performance metrics.}
\TODO{Table of status, exploitation, and yield metrics.}

\begin{table}[H]
\centering
\caption{Celtic Sea cod --- illustrative metrics table (replace with results).}
\label{tab:cod-metrics}
\begin{tabular}{@{}lrrr@{}}
\toprule
Metric & Realised & Counterfactual & Difference \\
\midrule
Years SSB $< B_{\lim}$ & --- & --- & --- \\
Years SSB $< MSYB_{\mathrm{trigger}}$ & --- & --- & --- \\
Mean $F/F_{\mathrm{MSY}}$ & --- & --- & --- \\
Mean catch & --- & --- & --- \\
\bottomrule
\end{tabular}
\end{table}

\subsection{Celtic Sea whiting (\texttt{whg.27.7b-ce-k})}
\label{sec:whg.27.7b-ce-k}
\TODO{Fill following the same pattern as Section~\ref{sec:cod.27.7e-k}.}

\subsection{Celtic Sea haddock (\texttt{had.27.7b-k})}
\label{sec:had.27.7b-k}
\TODO{Fill following the same pattern as Section~\ref{sec:cod.27.7e-k}.}

\subsection{Irish Sea whiting (\texttt{whg.27.7a})}
\label{sec:whg.27.7a}
\TODO{Fill following the same pattern as Section~\ref{sec:cod.27.7e-k}.}

\subsection{Irish Sea cod (\texttt{cod.27.7a})}
\label{sec:cod.27.7a}
\TODO{Fill following the same pattern as Section~\ref{sec:cod.27.7e-k}.}

\subsection{Northeast Atlantic mackerel (\texttt{mac.27.nea})}
\label{sec:mac.27.nea}
\TODO{Fill following the same pattern as Section~\ref{sec:cod.27.7e-k};
note Coastal States rule specifics.}

\FloatBarrier

%------------------------------------------------------------------------------
\section{Cross-stock synthesis}
\label{sec:synthesis}

\TODO{Where adherence would have improved outcomes; where effect was limited;
where shortcomings reflect the advice rule or assessment revision rather than
implementation failure.}

\placeholderfig{Cross-stock comparison of realised vs counterfactual status}{figs/synthesis_status.pdf}

\begin{table}[H]
\centering
\caption{Cross-stock summary (placeholder).}
\label{tab:synthesis}
\begin{tabular}{@{}lccc@{}}
\toprule
Stock & Status improved? & Yield change & Main driver \\
\midrule
Celtic Sea cod & --- & --- & --- \\
Celtic Sea whiting & --- & --- & --- \\
Celtic Sea haddock & --- & --- & --- \\
Irish Sea whiting & --- & --- & --- \\
Irish Sea cod & --- & --- & --- \\
NE Atlantic mackerel & --- & --- & --- \\
\bottomrule
\end{tabular}
\end{table}

%------------------------------------------------------------------------------
\section{Discussion}
\label{sec:discussion}

\subsection{Perception, recruitment, and reference points}

\TODO{Expand using week~1 discussion: (i)~perception changes at updates /
benchmarks; (ii)~recruitment impact; (iii)~reference points change at
benchmarks---and how the backtest design addresses each.}

\subsection{Limitations}

\TODO{Assessment approximation error; advice-rule simplifications; recruitment
assumptions; data gaps (e.g.\ Irish Sea whiting); perfect-implementation
idealisation.}

\subsection{Implications for advice and management}

\TODO{HCR design, rebuilding, and value of following advice vs revising
assessments.}

%------------------------------------------------------------------------------
\section{Conclusions}
\label{sec:conclusions}

\TODO{Bullet conclusions tied to objectives; confirmed stocks; headline
counterfactual findings; recommended follow-up.}

%------------------------------------------------------------------------------
\section{Software and reproducibility}
\label{sec:software}

Analysis uses FLR (\texttt{FLCore}, \texttt{FLBRP}, \texttt{FLasher},
\texttt{ggplotFL}) and \texttt{FLBacktest} / \texttt{backtest}
(\texttt{hcrICES}, peel and projection helpers). Notebooks write
\texttt{data/results/*.RData}; this report loads figures/tables generated from
those objects.

\begin{table}[H]
\centering
\caption{Key paths.}
\label{tab:paths}
\begin{tabular}{@{}ll@{}}
\toprule
Path & Role \\
\midrule
\texttt{Rmd/01\_screening.Rmd} & Task~1 screening \\
\texttt{Rmd/02.0\_condition\_om.Rmd} & OM conditioning (dynamic-\(B_0\) SRR) \\
\texttt{Rmd/02.1\_lterm\_eq.Rmd} & Long-term open-loop \(F_{\mathrm{MSY}}\) \\
\texttt{Rmd/03.0\_assessment.Rmd} & Task~2 assessment reconstruction \\
\texttt{Rmd/04.1\_openLoop.Rmd} / \texttt{04.2\_closedLoop.Rmd} & Task~3 counterfactuals \\
\texttt{Rmd/04.3\_backtest.Rmd} & Task~3 index \\
\texttt{Rmd/05.1\_metrics.Rmd} & Task~4 evaluation metrics \\
\texttt{Rmd/06.0\_report.Rmd} & Task~4 technical report (follows project-markup) \\
\texttt{scripts/run\_pipeline.R} & Knit pipeline in workplan order \\
\texttt{scripts/build\_week1\_tex.R} & Screening LaTeX tables/figures \\
\texttt{tex/backtest\_screening.tex} & Screening deliverable \\
\texttt{tex/project-markup.tex} & Contract workplan / report outline \\
\texttt{tex/report\_template.tex} & This technical report template \\
\texttt{data/results/01\_screening.RData} & Task~1 screening list \\
\texttt{data/results/02\_assessment.RData} & Task~2 assessment reconstruction \\
\texttt{data/results/03.0\_om.RData} & Conditioned OMs + realised summaries \\
\texttt{data/results/03.1\_openLoop.RData} & Open-loop benchmark projections \\
\texttt{data/results/03.2\_closedLoop.RData} & Closed-loop HCR counterfactuals \\
\texttt{R/paths.R} & Canonical results path helpers \\
\bottomrule
\end{tabular}
\end{table}

\paragraph{Compile.}
\begin{verbatim}
cd tex
pdflatex report_template.tex
pdflatex report_template.tex
\end{verbatim}

%------------------------------------------------------------------------------
\section*{Acknowledgements}
\TODO{Optional.}

\appendix
\section{Supplementary tables and figures}
\label{sec:appendix}
\TODO{Extra diagnostics, peel-by-peel plots, sensitivity runs.}

\end{document}
